\scanpage{12}
\mspage{10}
Noter que si $F$ est compatible avec $T^\ell_\xi$, il est à fortiori \uline{associé à} $\xi$, et on trouve
\[
G\deleted{\unsure{(Q_\ell)}}\longrightarrow\pi_1(S,\xi)
\]
d'où
\[
\unsure{\eta|\gamma}_\ell:G(Q_\ell)\longrightarrow\pi_1(S,\xi).
\]
On voit tout de suite que cet homomorphisme est inverse à gauche de $\unsure{\rho|\varphi}_\ell$
\[
\boxed{\eta_\ell\rho_\ell=\mathrm{id}_{\pi_1(S,\xi)}}
\]
Si on s'intéresse à une sous-catégorie $C_\alpha$ de $C$, correspondant à un $G_\alpha^{\unsure{0}}$, alors $G_\alpha/G_\alpha^0$ \deleted{correspond à}\inserted{définit} un quotient $\pi_\alpha$ de $\pi_1(S,\xi)$, et on a
\diagram{\begin{tikzcd}[ampersand replacement=\&,column sep=small,row sep=small]
\& \& G_\alpha/G_\alpha^0\arrow[dr,dashed] \& \\
\pi_1(S,\xi)\arrow[r,"\rho_\ell"]\arrow[rrr,bend right=27,"\text{\uline{hom can.}} "'] \& G_\alpha(Q_\ell)\arrow[ur,dashed]\arrow[rr] \& \& \pi_\alpha
\end{tikzcd}\quad(commutatif).}
La conjecture de Tate \inserted{(avec action semi-simple de Galois !)} équivaut à dire que, si $S$ est localisé d'un schéma de type fini sur $\mathrm Z$, alors $G_\alpha\bigl(\unsure{\to|-}Q_\ell\bigr)$ est l'enveloppe algébrique de $\operatorname{Im}(\pi_1(S,\xi)\unsure{-|\to}G_{\unsure{\ell|\alpha}}(Q_\ell))$.
% slanted
\margin{\uline{NB} La conjecture de Tate sous la forme : les $\frac12$-invariants dans $H^{2i}(\bar X,Q_\ell(i))$ sont les comb. à coeff. dans $Q_\ell$ des classes algébriques -- implique (avec conjectures standard) la $\frac12$ simplicité dans Tate\ldots}
\margin{\deleted{NB \unsure{il s'ensuit|$G_\alpha(Q_\ell)$} \illegible{3} \unsure{que l'action} \unsure{sur} $T_\ell(M)$, \unsure{algébriquement} \unsure{engendrée}}}

\uline{Cas} $S=\operatorname{Spec} k$, $k$ \uline{corps fini}.

Soit $C\subset M(k)$, et $F$ un ``foncteur fibre sur $C$ à coeff dans $K$. Comme on a un élément canonique \deleted{de}
\[
\mathrm{Frob}_k\in\operatorname{Aut}\underline{\mathrm{Id}}_{M(k)}\longrightarrow\operatorname{Aut}\mathrm{id}_C
\left[\begin{gathered}\simeq\unsure{Z|\mathcal Z|\mathfrak Z}(Q)\\
\text{où }Z\text{ est le centre de }G\\
\deleted{\unsure{C/F}}
\end{gathered}\right]
\]
on en déduit un élément canonique de $G(K)=\operatorname{Aut}F$, et même de $Z(Q)$ (où $Z$ est le centre de $G$)
\[
\boxed{\unsure{\xi|*}\ \mathrm{Frob}_{k\unsure{,|/}F}\in Z(Q)\subset G(K)}.
\]
\margin{N.B. $Z$ est en fait défini sur $Q$ !}
% archive sheet number: 11

